\Chapter{19}

\Verse{1}
{
    \zA{话说贾妃回宫，次日见驾谢恩，并回奏归省之事。 龙颜甚悦，又发内帑彩缎金 银等物以赐贾政及各椒房等员，不必细说。
        且说荣宁二府中连日用尽心力，真是人人力倦，各各神疲，又将园中一应陈设 动用之物，收拾了两三天方完。 第一个凤姐事多任重，别人或可偷闲躲静，独他是
        不能脱得的； 二则本性要强，不肯落人褒贬，只扎挣着与无事的人一样。 第一个宝 玉是极无事最闲暇的。
        偏这一早，袭人的母亲又亲来回过贾母，接袭人家去吃年茶， 晚上才得回来。}
}
{
    \eA{The Chia consort, we must now go on to explain, returned to the Palace, and
        the next day, on her appearance in the presence of His Majesty, she thanked
        him for his bounty and gave him furthermore an account of her experiences on
        her visit home.}
}
{
}

\Verse{2}
{
    \zA{因此，宝玉只和众丫头们掷骰子赶围棋作戏。 正在房内玩得没兴头， 忽见丫头们来回说：“东府里珍大爷来请过去看戏，放花灯。” 宝玉听了，便命换 衣裳。
        才要去时，忽又有贾妃赐出糖蒸酥酪来。 宝玉想上次袭人喜吃此物，便命留 与袭人了，自己回过贾母，过去看戏。
        谁想贾珍这边唱的是《丁郎认父》、《黄伯央大摆阴魂阵》，更有《孙行者大 闹天宫》、《姜太公斩将封神》等类的戏文。 倏尔神鬼乱出，忽又妖魔毕露。}
}
{
    \eA{His Majesty's dragon countenance was much elated, and he also issued from
        the privy store coloured satins, gold and silver and such like articles to
        be presented to Chia Cheng and the other officials in the various households
        of her relatives. But dispensing with minute details about them, we will now
        revert to the two mansions of Jung and Ning.}
}
{
}

\Verse{3}
{
    \zA{内中 扬幡过会、号佛行香、锣鼓喊叫之声，闻于巷外。 弟兄子侄，互为献酬； 姊妹婢妾， 共相笑语。
        独有宝玉见那繁华热闹到如此不堪的田地，只略坐了一坐，便走往各处 闲耍。 先是进内去和尤氏并丫头姬妾鬼混了一回，便出二门来。 尤氏等仍料他出来
        看戏，遂也不曾照管。 贾珍、贾琏、薛蟠等只顾猜谜行令，百般作乐，纵一时不见 他在座，只道在里边去了，也不理论。}
}
{
    \eA{With the extreme strain on mind and body for successive days, the strength
        of one and all was, in point of fact, worn out and their respective energies
        exhausted. And it was besides after they had been putting by the various
        decorations and articles of use for two or three days, that they, at length,
        got through the work. Lady Feng was the one who had most to do, and whose
        responsibilities were greatest.}
}
{
}

\Verse{4}
{
    \zA{至于跟宝玉的小厮们，那年纪大些的，知宝 玉这一来了必是晚上才散，因此偷空儿也有会赌钱的，也有往亲友家去的，或赌或 饮，都私自散了，待晚上再来；
        那些小些的，都钻进戏房里瞧热闹儿去了。 宝玉见一个人没有，因想：“素日这里有个小书房内曾挂着一轴美人，画的很 得神。
        今日这般热闹，想那里自然无人，那美人也自然是寂寞的，须得我去望慰他 一回。” 想着，便往那里来。 刚到窗前，听见屋里一片喘息之声。 宝玉倒唬了一跳，
        心想：“美人活了不成？”}
}
{
    \eA{The others could possibly steal a few leisure moments and retire to rest,
        while she was the sole person who could not slip away. In the second place,
        naturally anxious as she was to excel and both to fall in people's
        estimation, she put up with the strain just as if she were like one of those
        who had nothing to attend to. But the one who had the least to do and had
        the most leisure was Pao-yü.}
}
{
}

\Verse{5}
{
    \zA{乃大着胆子，舐破窗纸。 向内一看，那轴美人却不曾活， 却是茗烟按着个女孩子，也干那警幻所训之事，正在得趣，故此呻吟。
        宝玉禁不住，大叫“了不得”，一脚踹进门去。 将两个唬的抖衣而颤。 茗烟见 是宝玉，忙跪下哀求。 宝玉道：“青天白日，这是怎么说!珍大爷要知道了，你是
        死是活？” 一面看那丫头，倒也白白净净儿的有些动人心处，在那里羞的脸红耳赤， 低首无言。 宝玉跺脚道：“还不快跑！” 一语提醒，那丫头飞跑去了。}
}
{
    \eA{As luck would have it on this day, at an early hour, Hsi Jen's mother came
        again in person and told dowager lady Chia that she would take Hsi Jen home
        to drink a cup of tea brewed in the new year and that she would return in
        the evening. For this reason Pao-yü was only in the company of all the
        waiting-maids, throwing dice, playing at chess and amusing himself.}
}
{
}

\Verse{6}
{
    \zA{宝玉又赶出 去叫道：“你别怕，我不告诉人！” 急的茗烟在后叫：“祖宗，这是分明告诉人了！” 宝玉因问：“那丫头十几岁了？” 茗烟道：“不过十六七了。”
        宝玉道：“连他的 岁数也不问问，就作这个事，可见他白认得你了。 可怜，可怜！” 又问：“名字叫 什么？” 茗烟笑道：“若说出名字来话长，真正新鲜奇文。
        他说他母亲养他的时节， 做了一个梦，梦得了一匹锦，上面是五色富贵不断头的‘？’ 字花样，所以他的名 字就叫做万儿。”}
}
{
    \eA{But while he was in the room playing with them with a total absence of zest,
        he unawares perceived a few waiting-maids arrive, who informed him that
        their senior master Mr. Chen, of the Eastern Mansion, had come to invite him
        to go and see a theatrical performance, and the fireworks, which were to be
        let off. Upon hearing these words, Pao-yü speedily asked them to change his
        clothes;}
}
{
}

\Verse{7}
{
    \zA{宝玉听了笑道：“想必他将来有些造化。 等我明儿说了给你作媳 妇，好不好？” 茗烟也笑了。 因问：“二爷为何不看这样的好戏？” 宝玉道：“看
        了半日，怪烦的，出来逛逛，就遇见你们了。 这会子作什么呢？” 茗烟微微笑道： “这会子没人知道，我悄悄的引二爷城外逛去，一会儿再回这里来。”
        宝玉道：“不 好，看仔细花子拐了去。 况且他们知道了，又闹大了。 不如往近些的地方去，还可 就来。” 茗烟道：“就近地方谁家可去？ 这却难了。”}
}
{
    \eA{but just as he was ready to start, presents of cream, steamed with sugar,
        arrived again when least expected from the Chia Consort, and Pao-yü
        recollecting with what relish Hsi Jen had partaken of this dish on the last
        occasion forthwith bid them keep it for her; while he went himself and told
        dowager lady Chia that he was going over to see the play.}
}
{
}

\Verse{8}
{
    \zA{宝玉笑道：“依我的主意， 咱们竟找花大姐姐去，瞧他在家作什么呢。” 茗烟笑道：“好!好!倒忘了他家。”
        又道：“他们知道了，说我引着二爷胡走，要打我呢。” 宝玉道：“有我呢！” 茗 烟听说，拉了马，二人从后门就走了。
        幸而袭人家不远，不过一半里路程，转眼已到门前。 茗烟先进去叫袭人之兄花 自芳。}
}
{
    \eA{The plays sung over at Chia Chen's consisted, who would have thought it, of
        "Ting L'ang recognises his father," and "Huang Po-ying deploys the spirits
        for battle," and in addition to these, "Sung Hsing-che causes great
        commotion in the heavenly palace;" "Ghiang T'ai-kung kills the general and
        deifies him," and other such like.}
}
{
}

\Verse{9}
{
    \zA{此时袭人之母接了袭人与几个外甥女儿几个侄女儿来家，正吃果茶，听见外 面有人叫“花大哥”，花自芳忙出去看时，见是他主仆两个，唬的惊疑不定，连忙
        抱下宝玉来，至院内嚷道：“宝二爷来了！” 别人听见还可，袭人听了，也不知为 何，忙跑出来迎着宝玉，一把拉着问：“你怎么来了？” 宝玉笑道：“我怪闷的，
        来瞧瞧你作什么呢。” 袭人听了，才把心放下来，说道：“你也胡闹了!可作什么 来呢？” 一面又问茗烟：“还有谁跟了来了？”}
}
{
    \eA{Soon appeared the spirits and devils in a confused crowd on the stage, and
        suddenly also became visible the whole band of sprites and goblins, among
        which were some waving streamers, as they went past in a procession,
        invoking Buddha and burning incense. The sound of the gongs and drums and of
        shouts and cries were audible at a distance beyond the lane;}
}
{
}

\Verse{10}
{
    \zA{茗烟笑道：“别人都不知道。” 袭 人听了，复又惊慌道：“这还了得!倘或碰见人，或是遇见老爷，街上人挤马碰， 有个失闪，这也是玩得的吗？
        你们的胆子比斗还大呢!都是茗烟调唆的，等我回去告 诉嬷嬷们，一定打你个贼死。” 茗烟撅了嘴道：“爷骂着打着叫我带了来的，这会 子推到我身上。
        我说别来罢!――要不，我们回去罢。” 花自芳忙劝道：“罢了， 已经来了，也不用多说了。 只是茅檐草舍，又窄又不干净，爷怎么坐呢？”
        袭人的母亲也早迎出来了。}
}
{
    \eA{and in the whole street, one and all extolled the performance as
        exceptionally grand, and that the like could never have been had in the
        house of any other family. Pao-yü, noticing that the commotion and bustle
        had reached a stage so unbearable to his taste, speedily betook himself,
        after merely sitting for a little while, to other places in search of
        relaxation and fun.}
}
{
}

\Verse{11}
{
    \zA{袭人拉着宝玉进去。 宝玉见房中三五个女孩儿，见 他进来，都低了头，羞的脸上通红。 花自芳母子两个恐怕宝玉冷，又让他上炕，又 忙另摆果子，又忙倒好茶。
        袭人笑道：“你们不用白忙，我自然知道，不敢乱给他 东西吃的。” 一面说，一面将自己的坐褥拿了来，铺在一个杌子上，扶着宝玉坐下，
        又用自己的脚炉垫了脚，向荷包内取出两个梅花香饼儿来，又将自己的手炉掀开焚 上，仍盖好，放在宝玉怀里，然后将自己的茶杯斟了茶，送与宝玉。}
}
{
    \eA{First of all, he entered the inner rooms, and after spending some time in
        chatting and laughing with Mrs. Yu, the waiting-maids, and secondary wives,
        he eventually took his departure out of the second gate; and as Mrs. Yu and
        her companions were still under the impression that he was going out again
        to see the play, they let him speed on his way, without so much as keeping
        an eye over him.}
}
{
}

\Verse{12}
{
    \zA{彼时他母兄已 是忙着齐齐整整的摆上一桌子果品来，袭人见总无可吃之物，因笑道：“既来了， 没有空回去的理，好歹尝一点儿，也是来我家一趟。”
        说着，捻了几个松瓤，吹去 细皮，用手帕托着给他。 宝玉看见袭人两眼微红，粉光融滑，因悄问袭人道：“好好的哭什么？” 袭人 笑道：“谁哭来着？
        才迷了眼揉的。” 因此便遮掩过了。 因见宝玉穿着大红金蟒狐 腋箭袖，外罩石青貂裘排穗褂，说道：“你特为往这里来，又换新衣裳，他们就不 问你往那里去吗？”}
}
{
    \eA{Chia Chen, Chia Lien, Hsúeh P'an and the others were bent upon guessing
        enigmas, enforcing the penalties and enjoying themselves in a hundred and
        one ways, so that even allowing that they had for a moment noticed that he
        was not occupying his seat, they must merely have imagined that he had gone
        inside and not, in fact, worried their minds about him.}
}
{
}

\Verse{13}
{
    \zA{宝玉道：“原是珍大爷请过去看戏换的。” 袭人点头，又道： “坐一坐就回去罢，这个地方儿不是你来得的。” 宝玉笑道：“你就家去才好呢，
        我还替你留着好东西呢。” 袭人笑道：“悄悄儿的罢!叫他们听着作什么？” 一面 又伸手从宝玉项上将通灵玉摘下来，向他姊妹们笑道：“你们见识见识。 时常说起
        来都当稀罕，恨不能一见，今儿可尽力儿瞧瞧。 再瞧什么稀罕物儿，也不过是这么 着了。” 说毕递与他们，传看了一遍，仍与宝玉挂好。}
}
{
    \eA{And as for the pages, who had come along with Pao-yü, those who were a
        little advanced in years, knowing very well that Pao-yü would, on an
        occasion like the present, be sure not to be going before dusk, stealthily
        therefore took advantage of his absence, those, who could, to gamble for
        money, and others to go to the houses of relatives and friends to drink of
        the new year tea, so that what with gambling and drinking the whole bevy
        surreptitiously dispersed, waiting for dusk before they came back;}
}
{
}

\Verse{14}
{
    \zA{又命他哥哥去雇一辆干干净 净、严严紧紧的车，送宝玉回去。 花自芳道：“有我送去，骑马也不妨了。” 袭人 道：“不为不妨，为的是碰见人。”
        花自芳忙去雇了一辆车来，众人也不好相留， 只得送宝玉出去。 袭人又抓些果子给茗烟，又把些钱给他买花爆放，叫他：“别告 诉人，连你也有不是。”
        一面说着，一直送宝玉至门前，看着上车，放下车帘。 茗 烟二人牵马跟随。}
}
{
    \eA{while those, who were younger, had all crept into the green rooms to watch
        the excitement; with the result that Pao-yü perceiving not one of them about
        bethought himself of a small reading room, which existed in previous days on
        this side, in which was suspended a picture of a beauty so artistically
        executed as to look life-like.}
}
{
}

\Verse{15}
{
    \zA{来至宁府街，茗烟命住车，向花自芳道：“须得我和二爷还到东 府里混一混，才过去得呢，看大家疑惑。” 花自芳听说有理，忙将宝玉抱下车来， 送上马去。
        宝玉笑说：“倒难为你了。” 于是仍进了后门来，俱不在话下。 却说宝玉自出了门，他房中这些丫鬟们都索性恣意的玩笑，也有赶围棋的，也
        有掷骰抹牌的，磕了一地的瓜子皮儿，偏奶母李嬷嬷拄拐进来请安，瞧瞧宝玉； 见 宝玉不在家，丫鬟们只顾玩闹，十分看不过。}
}
{
    \eA{"On such a bustling day as this," he reasoned, "it's pretty certain, I
        fancy, that there will be no one in there; and that beautiful person must
        surely too feel lonely, so that it's only right that I should go and console
        her a bit." With these thoughts, he hastily betook himself towards the
        side-house yonder, and as soon as he came up to the window, he heard the
        sound of groans in the room.}
}
{
}

\Verse{16}
{
    \zA{因叹道：“只从我出去了不大进来， 你们越发没了样儿了，别的嬷嬷越不敢说你们了。 那宝玉是个‘丈八的灯台――照 见人家，照不见自己’的，只知嫌人家腌？。
        这是他的房子，由着你们遭塌，越不 成体统了。” 这些丫头们明知宝玉不讲究这些，二则李嬷嬷已是告老解事出去的了， 如今管不着他们。
        因此，只顾玩笑，并不理他。 那李嬷嬷还只管问：“宝玉如今一 顿吃多少饭？ 什么时候睡觉？” 丫头们总胡乱答应，有的说：“好个讨厌的老货！”}
}
{
    \eA{Pao-yü was really quite startled. "What!" (he thought), "can that beautiful
        girl, possibly, have come to life!" and screwing up his courage, he licked a
        hole in the paper of the window and peeped in. It was not she, however, who
        had come to life, but Ming Yen holding down a girl and likewise indulging in
        what the Monitory Dream Fairy had taught him. "Dreadful!"}
}
{
}

\Verse{17}
{
    \zA{李嬷嬷又问道：“这盖碗里是酪，怎么不送给我吃？” 说毕，拿起就吃。 一个 丫头道：“快别动!那是说了给袭人留着的，回来又惹气了。 你老人家自己承认，
        别带累我们受气。” 李嬷嬷听了，又气又愧，便说道：“我不信他这么坏了肠子! 别说我吃了一碗牛奶，就是再比这个值钱的，也是应该的。 难道待袭人比我还重？
        难道他不想想怎么长大了？ 我的血变了奶，吃的长这么大，如今我吃他碗牛奶，他 就生气了？}
}
{
    \eA{exclaimed Pao-yü, aloud, unable to repress himself, and, stamping one of his
        feet, he walked into the door to the terror of both of them, who parting
        company, shivered with fear, like clothes that are being shaken. Ming Yen
        perceiving that it was Pao-yü promptly fell on his knees and piteously
        implored for pardon. "What! in broad daylight! what do you mean by it?}
}
{
}

\Verse{18}
{
    \zA{我偏吃了，看他怎么着!你们看袭人不知怎么样，那是我手里调理出来的 毛丫头，什么阿物儿！” 一面说，一面赌气把酪全吃了。 又一个丫头笑道：“他们
        不会说话，怨不得你老人家生气。 宝玉还送东西给你老人家去，岂有为这个不自在 的？” 李嬷嬷道：“你也不必妆狐媚子哄我，打量上次为茶撵茜雪的事我不知道呢!
        明儿有了不是，我再来领。” 说着，赌气去了。 少时，宝玉回来，命人去接袭人，只见晴雯躺在床上不动，宝玉因问：“可是 病了？ 还是输了呢？”}
}
{
    \eA{Were your master Mr. Chen to hear of it, would you die or live?" asked
        Pao-yü, as he simultaneously cast a glance at the servant-girl, who although
        not a beauty was anyhow so spick and span, and possessed besides a few
        charms sufficient to touch the heart. From shame, her face was red and her
        ears purple, while she lowered her head and uttered not a syllable. Pao-yü
        stamped his foot. "What!"}
}
{
}

\Verse{19}
{
    \zA{秋纹道：“他倒是赢的； 谁知李老太太来了混输了，他气的 睡去了。” 宝玉笑道：“你们别和他一般见识，由他去就是了。” 说着，袭人已来，彼此相见。
        袭人又问宝玉何处吃饭，多早晚回来； 又代母妹 问诸同伴姊妹好。 一时换衣卸妆。 宝玉命取酥酪来，丫鬟们回说：“李奶奶吃了。”
        宝玉才要说话，袭人便忙笑说道：“原来留的是这个，多谢费心。 前儿我因为好吃， 吃多了，好肚子疼，闹的吐了才好了。 他吃了倒好，搁在这里白遭塌了。}
}
{
    \eA{he shouted, "don't you yet bundle yourself away!" This simple remark
        suggested the idea to the girl's mind who ran off, as if she had wings to
        fly with; but as Pao-yü went also so far as to go in pursuit of her, calling
        out: "Don't be afraid, I'm not one to tell anyone," Ming Yen was so
        exasperated that he cried, as he went after them, "My worthy ancestor, this
        is distinctly telling people about it."}
}
{
}

\Verse{20}
{
    \zA{我只想风 干栗子吃，你替我剥栗子，我去铺炕。” 宝玉听了，信以为真，方把酥酪丢开，取 了栗子来，自向灯下检剥。
        一面见众人不在房中，乃笑问袭人道：“今儿那个穿红 的是你什么人？” 袭人道：“那是我两姨姐姐。” 宝玉听了，赞叹了两声。 袭人道： “叹什么？
        我知道你心里的缘故。 想是说：他那里配穿红的？” 宝玉笑道：“不是 不是。 那样的人不配穿红的，谁还敢穿？ 我因为见他实在好的很，怎么也得他在咱
        们家就好了。”}
}
{
    \eA{"How old is that servant girl?" Pao-yü having asked; "She's, I expect, no
        more than sixteen or seventeen," Ming Yen rejoined. "Well, if you haven't
        gone so far as to even ascertain her age," Pao-yü observed, "you're sure to
        know still less about other things; and it makes it plain enough that her
        acquaintance with you is all vain and futile! What a pity! what a pity!" He
        then went on to enquire what her name was;}
}
{
}

\Verse{21}
{
    \zA{袭人冷笑道：“我一个人是奴才命罢了，难道连我的亲戚都是奴才 命不成？ 定还要拣实在好的丫头才往你们家来？” 宝玉听了，忙笑道：“你又多心
        了!我说往咱们家来，必定是奴才不成，说亲戚就使不得？” 袭人道：“那也搬配 不上。” 宝玉便不肯再说，只是剥栗子。 袭人笑道：“怎么不言语了？
        想是我才冒撞冲 犯了你？ 明儿赌气花几两银子买进他们来就是了。” 宝玉笑道：“你说的话怎么叫 人答言呢？}
}
{
    \eA{and "Were I," continued Ming Yen smiling, "to tell you about her name it
        would involve a long yarn; it's indeed a novel and strange story! She
        relates that while her mother was nursing her, she dreamt a dream and
        obtained in this dream possession of a piece of brocaded silk, on which were
        designs, in variegated colours, representing opulence and honour, and a
        continuous line of the character Wan;}
}
{
}

\Verse{22}
{
    \zA{我不过是赞他好，正配生在这深宅大院里，没的我们这宗浊物倒生在这 里！” 袭人道：“他虽没这样造化，倒也是娇生惯养的，我姨父姨娘的宝贝儿似的，
        如今十七岁，各样的嫁妆都齐备了，明年就出嫁。” 宝玉听了“出嫁”二字，不禁 又？ 了两声。 正不自在，又听袭人叹道：“我这几年，姊妹们都不大见。
        如今我要 回去了，他们又都去了！” 宝玉听这话里有文章，不觉吃了一惊，忙扔下栗子，问 道：“怎么着，你如今要回去？”}
}
{
    \eA{and that this reason accounts for the name of Wan Erh, which was given her."
        "This is really strange!" Pao-yü exclaimed with a grin, after lending an ear
        to what he had to say; "and she is bound, I think, by and by to have a good
        deal of good fortune!"}
}
{
}

\Verse{23}
{
    \zA{袭人道：“我今儿听见我妈和哥哥商量，教我再 耐一年，明年他们上来就赎出我去呢。” 宝玉听了这话，越发忙了，因问：“为什 么赎你呢？”
        袭人道：“这话奇了!我又比不得是这里的家生子儿，我们一家子都 在别处，独我一个人在这里，怎么是个了手呢？” 宝玉道：“我不叫你去也难哪！”
        袭人道：“从来没这个理。 就是朝廷宫里，也有定例，几年一挑，几年一放，没有 长远留下人的理，别说你们家！”
        宝玉想一想，果然有理，又道：“老太太要不放你呢？”}
}
{
    \eA{These words uttered, he plunged in deep thought for a while, and Ming Yen
        having felt constrained to inquire: "Why aren't you, Mr. Secundus, watching
        a theatrical performance of this excellent kind?" "I had been looking on for
        ever so long," Pao-yü replied, "until I got quite weary; and had just come
        out for a stroll, when I happened to meet you two. But what's to be done
        now?" Ming Yen gave a faint smile.}
}
{
}

\Verse{24}
{
    \zA{袭人道：“为什么不 放呢？ 我果然是个难得的，或者感动了老太太、太太不肯放我出去，再多给我们家 几两银子留下，也还有的；
        其实我又不过是个最平常的人，比我强的多而且多。 我 从小儿跟着老太太，先伏侍了史大姑娘几年，这会子又伏侍了你几年，我们家要来
        赎我，正是该叫去的，只怕连身价不要就开恩放我去呢。 要说为伏侍的你好不叫我 去，断然没有的事。 那伏侍的好，是分内应当的，不是什么奇功； 我去了仍旧又有
        好的了，不是没了我就使不得的。”}
}
{
    \eA{"As there's no one here to know anything about it," he added, "I'll
        stealthily take you, Mr. Secundus, for a walk outside the city walls; and
        we'll come back shortly, before they've got wind of it." "That won't do,"
        Pao-yü demurred, "we must be careful, or else some beggar might kidnap us
        away; besides, were they to come to hear of it, there'll be again a dreadful
        row;}
}
{
}

\Verse{25}
{
    \zA{宝玉听了这些话，竟是有去的理无留的理，心 里越发急了，因又道：“虽然如此说，我的一心要留下你，不怕老太太不和你母亲
        说，多多给你母亲些银子，他也不好意思接你了。” 袭人道：“我妈自然不敢强。 且慢说和他好说，又多给银子； 就便不好和他说，一个钱也不给，安心要强留下我，
        他也不敢不依。 但只是咱们家从没干过这倚势仗贵霸道的事。 这比不得别的东西， 因为喜欢，加十倍利弄了来给你，那卖的人不吃亏，就可以行得的；}
}
{
    \eA{and isn't it better that we should go to some nearer place, from which we
        could, after all, return at once?" "As for some nearer place," Ming Yen
        observed; "to whose house can we go? It's really no easy matter!" "My idea
        is," Pao-yü suggested with a smirk, "that we should simply go, and find
        sister Hua, and see what she's up to at home." "Yes! Yes!" Ming Yen replied
        laughingly;}
}
{
}

\Verse{26}
{
    \zA{如今无故平空 留下我，于你又无益，反教我们骨肉分离，这件事，老太太、太太肯行吗？” 宝玉 听了，思忖半晌，乃说道：“依你说来说去，是去定了？”
        袭人道：“去定了。” 宝玉听了自思道：“谁知这样一个人，这样薄情无义呢！” 乃叹道：“早知道都是 要去的，我就不该弄了来。 临了剩我一个孤鬼儿！”
        说着便赌气上床睡了。 原来袭人在家，听见他母兄要赎他回去，他就说：“至死也不回去。”}
}
{
    \eA{"the fact is I had forgotten all about her home; but should it reach their
        ears," he continued, "they'll say that it was I who led you, Mr. Secundus,
        astray, and they'll beat me!" "I'm here for you!" Pao-yü having assured him;
        Ming Yen at these words led the horses round, and the two of them speedily
        made their exit by the back gate. Luckily Hsi Jen's house was not far off.}
}
{
}

\Verse{27}
{
    \zA{又说： “当日原是你们没饭吃，就剩了我还值几两银子，要不叫你们卖，没有个看着老子 娘饿死的理； 如今幸而卖到这个地方儿，吃穿和主子一样，又不朝打暮骂。
        况如今 爹虽没了，你们却又整理的家成业就，复了元气。 若果然还艰难，把我赎出来再多 掏摸几个钱，也还罢了，其实又不难了。 这会子又赎我做什么？
        权当我死了，再不 必起赎我的念头了！” 因此哭了一阵。 他母兄见他这般坚执，自然必不出来的了。}
}
{
    \eA{It was no further than half a li's distance, so that in a twinkle they had
        already reached the front of the door, and Ming Yen was the first to walk in
        and to call for Hsi Jen's eldest brother Hua Tzu-fang.}
}
{
}

\Verse{28}
{
    \zA{况且原是卖倒的死契，明仗着贾宅是慈善宽厚人家儿，不过求求，只怕连身价银一 并赏了还是有的事呢； 二则贾府中从不曾作践下人，只有恩多威少的，且凡老少房
        中所有亲侍的女孩子们，更比待家下众人不同，平常寒薄人家的女孩儿也不能那么 尊重：因此他母子两个就死心不赎了。 次后忽然宝玉去了，他两个又是那个光景儿，
        母子二人心中更明白了，越发一块石头落了地，而且是意外之想，彼此放心，再无 别意了。}
}
{
    \eA{Hsi Jen's mother had, on this occasion, united in her home Hsi Jen, several
        of her sister's daughters, as well as a few of her nieces, and they were
        engaged in partaking of fruits and tea, when they heard some one outside
        call out, "Brother Hua." Hua Tzu-fang lost no time in rushing out;}
}
{
}

\Verse{29}
{
    \zA{且说袭人自幼儿见宝玉性格异常，其淘气憨顽出于众小儿之外，更有几件千奇 百怪口不能言的毛病儿。 近来仗着祖母溺爱，父母亦不能十分严紧拘管，更觉放纵
        弛荡，任情恣性，最不喜务正。 每欲劝时，谅不能听。 今日可巧有赎身之论，故先 用骗词以探其情，以压其气，然后好下箴规。 今见宝玉默默睡去，知其情有不忍，
        气已馁堕。 自己原不想栗子吃，只因怕为酥酪生事，又像那茜雪之茶，是以假要栗 子为由，混过宝玉不提就完了。}
}
{
    \eA{and upon looking and finding that it was the two of them, the master and his
        servant, he was so taken by surprise that his fears could not be set at
        rest. Promptly, he clasped Pao-yü in his arms and dismounted him, and coming
        into the court, he shouted out at the top of his voice: "Mr. Pao has come."}
}
{
}

\Verse{30}
{
    \zA{于是命小丫头子们将栗子拿去吃了，自己来推宝玉。 只见宝玉泪痕满面，袭人便笑道：“这有什么伤心的？ 你果然留我，我自然不肯出 去。”
        宝玉见这话头儿活动了，便道：“你说说我还要怎么留你？ 我自己也难说了！” 袭人笑道：“咱们两个的好，是不用说了。 但你要安心留我，不在这上头。 我另说
        出三件事来，你果然依了，那就是真心留我了，刀搁在脖子上我也不出去了。” 宝玉忙笑道：“你说那几件？ 我都依你。}
}
{
    \eA{The other persons heard the announcement of his arrival, with equanimity,
        but when it reached Hsi Jen's ears, she truly felt at such a loss to fathom
        the object of his visit that issuing hastily out of the room, she came to
        meet Pao-yü, and as she laid hold of him: "Why did you come?" she asked. "I
        felt awfully dull," Pao-yü rejoined with a smile, "and came to see what you
        were up to."}
}
{
}

\Verse{31}
{
    \zA{好姐姐，好亲姐姐!别说两三件，就是 两三百件我也依的。 只求你们看守着我，等我有一日化成了飞灰，――飞灰还不好， 灰还有形有迹，还有知识的。
        ――等我化成一股轻烟，风一吹就散了的时候儿，你 们也管不得我，我也顾不得你们了，凭你们爱那里去那里去就完了。” 急的袭人忙
        握他的嘴，道：“好爷!我正为劝你这些个。 更说的狠了！” 宝玉忙说道：“再不 说这话了。” 袭人道：“这是头一件要改的。”
        宝玉道：“改了，再说你就拧嘴! 还有什么？”}
}
{
    \eA{Hsi Jen at these words banished, at last, all anxiety from her mind. "You're
        again up to your larks," she observed, "but what's the aim of your visit?
        Who else has come along with him?" she at the same time went on to question
        Ming Yen. "All the others know nothing about it!" explained Ming Yen
        exultingly; "only we two do, that's all."}
}
{
}

\Verse{32}
{
    \zA{袭人道：“第二件，你真爱念书也罢，假爱也罢，只在老爷跟前，或在别人跟 前，你别只管嘴里混批，只作出个爱念书的样儿来，也叫老爷少生点儿气，在人跟
        前也好说嘴。 老爷心里想着：我家代代念书，只从有了你，不承望不但不爱念书， 已经他心里又气又恼了，而且背前面后混批评。 凡读书上进的人，你就起个外号儿，
        叫人家‘禄蠹’； 又说只除了什么‘明明德’外就没书了，都是前人自己混编纂出 来的。 这些话你怎么怨得老爷不气，不时时刻刻的要打你呢？”}
}
{
    \eA{When Hsi Jen heard this remark, she gave way afresh to solicitous fears:
        "This is dreadful!" she added; "for were you to come across any one from the
        house, or to meet master; or were, in the streets, people to press against
        you, or horses to collide with you, as to make (his horse) shy, and he were
        to fall, would that too be a joke? The gall of both of you is larger than a
        peck measure;}
}
{
}

\Verse{33}
{
    \zA{宝玉笑道：“再不 说了。 那是我小时候儿不知天多高地多厚信口胡说的，如今再不敢说了。 还有什么 呢？” 袭人道：“再不许谤僧毁道的了。
        还有更要紧的一件事，再不许弄花儿，弄 粉儿，偷着吃人嘴上擦的胭脂，和那个爱红的毛病儿了。” 宝玉道：“都改!都改! 再有什么快说罢。”
        袭人道：“也没有了，只是百事检点些，不任意任性的就是了。 你要果然都依了，就拿八人轿也抬不出我去了。” 宝玉笑道：“你这里长远了，不 怕没八人轿你坐。”}
}
{
    \eA{but it's all you, Ming Yen, who has incited him, and when I go back, I'll
        surely tell the nurses to beat you." Ming Yen pouted his mouth. "Mr.
        Secundus," he pleaded, "abused me and beat me, as he bade me bring him here,
        and now he shoves the blame on my shoulders! 'Don't let us go,' I suggested;
        'but if you do insist, well then let us go and have done.'" Hua Tzu-fang
        promptly interceded. "Let things alone," he said;}
}
{
}

\Verse{34}
{
    \zA{袭人冷笑道：“这我可不希罕的。 有那个福气，没有那个道理， 纵坐了也没趣儿。” 二人正说着，只见秋纹走进来，说：“三更天了，该睡了。 方才老太太打发嬷
        嬷来问，我答应睡了。” 宝玉命取表来看时，果然针已指到子初二刻了，方从新盥 漱，宽衣安歇，不在话下。
        至次日清晨，袭人起来，便觉身体发重，头疼目胀，四肢火热。 先时还扎挣的 住，次后捱不住，只要睡，因而和衣躺在炕上。}
}
{
    \eA{"now that they're already here, there's no need whatever of much ado. The
        only thing is that our mean house with its thatched roof is both so crammed
        and so filthy that how could you, sir, sit in it!" Hsi Jen's mother also
        came out at an early period to receive him, and Hsi Jen pulled Pao-yü in.}
}
{
}

\Verse{35}
{
    \zA{宝玉忙回了贾母，传医诊视，说道： “不过偶感风寒，吃一两剂药疏散疏散就好了。” 开方去后，令人取药来煎好，刚
        服下去，命他盖上被窝渥汗，宝玉自去黛玉房中来看视。 彼时黛玉自在床上歇午，丫鬟们皆出去自便，满屋内静悄悄的。 宝玉揭起绣线
        软帘，进入里间，只见黛玉睡在那里，忙上来推他道：“好妹妹，才吃了饭，又睡 觉！” 将黛玉唤醒。 黛玉见是宝玉，因说道：“你且出去逛逛，我前儿闹了一夜，
        今儿还没歇过来，浑身酸疼。”}
}
{
    \eA{Once inside the room, Pao-yü perceived three or five girls, who, as soon as
        they caught sight of him approaching, all lowered their heads, and felt so
        bashful that their faces were suffused with blushes.}
}
{
}

\Verse{36}
{
    \zA{宝玉道：“酸疼事小，睡出来的病大，我替你解闷 儿，混过困去就好了。” 黛玉只合着眼，说道：“我不困，只略歇歇儿，你且别处 去闹会子再来。”
        宝玉推他道：“我往那里去呢，见了别人就怪腻的。” 黛玉听了， “嗤”的一笑道：“你既要在这里，那边去老老实实的坐着，咱们说话儿。” 宝玉
        道：“我也歪着。” 黛玉道：“你就歪着。” 宝玉道：“没有枕头。 咱们在一个枕 头上罢。” 黛玉道：“放屁!外头不是枕头？ 拿一个来枕着。”}
}
{
    \eA{But as both Hua Tzu-fang and his mother were afraid that Pao-yü would catch
        cold, they pressed him to take a seat on the stove-bed, and hastened to
        serve a fresh supply of refreshments, and to at once bring him a cup of good
        tea. "You needn't be flurrying all for nothing," Hsi Jen smilingly
        interposed; "I, naturally, should know;}
}
{
}

\Verse{37}
{
    \zA{宝玉出至外间，看了 一看，回来笑道：“那个我不要，也不知是那个腌？ 老婆子的。” 黛玉听了，睁开 眼，起身笑道：“真真你就是我命中的‘魔星’。
        请枕这一个！” 说着，将自己枕 的推给宝玉，又起身将自己的再拿了一个来枕上，二人对着脸儿躺下。
        黛玉一回眼，看见宝玉左边腮上有钮扣大小的一块血迹，便欠身凑近前来，以 手抚之细看道：“这又是谁的指甲划破了？” 宝玉倒身，一面躲，一面笑道：“不
        是划的，只怕是才刚替他们淘澄胭脂膏子溅上了一点儿。”}
}
{
    \eA{and there's no use of even laying out any fruits, as I daren't recklessly
        give him anything to eat." Saying this, she simultaneously took her own
        cushion and laid it on a stool, and after Pao-yü took a seat on it, she
        placed the footstove she had been using, under his feet; and producing, from
        a satchet, two peach-blossom-scented small cakes, she opened her own
        hand-stove and threw them into the fire;}
}
{
}

\Verse{38}
{
    \zA{说着，便找绢子要擦。 黛玉便用自己的绢子替他擦了，咂着嘴儿说道：“你又干这些事了。 干也罢了，必 定还要带出幌子来。
        就是舅舅看不见，别人看见了，又当作奇怪事新鲜话儿去学舌 讨好儿，吹到舅舅耳朵里，大家又该不得心净了。” 宝玉总没听见这些话，只闻见
        一股幽香，却是从黛玉袖中发出，闻之令人醉魂酥骨。 宝玉一把便将黛玉的衣袖拉 住，要瞧瞧笼着何物。 黛玉笑道：“这时候谁带什么香呢？”
        宝玉笑道：“那么着， 这香是那里来的？”}
}
{
    \eA{which done, she covered it well again and placed it in Pao-yü's lap. And
        eventually, she filled her own tea-cup with tea and presented it to Pao-yü,
        while, during this time, her mother and sister had been fussing about,
        laying out in fine array a tableful of every kind of eatables.}
}
{
}

\Verse{39}
{
    \zA{黛玉道：“连我也不知道，想必是柜子里头的香气熏染的，也 未可知。” 宝玉摇头道：“未必。 这香的气味奇怪，不是那些香饼子、香球子、香 袋儿的香。”
        黛玉冷笑道：“难道我也有什么‘罗汉’‘真人’给我些奇香不成？ 就是得了奇香，也没有亲哥哥亲兄弟弄了花儿、朵儿、霜儿、雪儿替我炮制。 我有
        的是那些俗香罢了！” 宝玉笑道：“凡我说一句，你就拉上这些。 不给你个利害也 不知道，从今儿可不饶你了！”}
}
{
    \eA{Hsi Jen noticed that there were absolutely no things that he could eat, but
        she felt urged to say with a smile: "Since you've come, it isn't right that
        you should go empty away; and you must, whether the things be good or bad,
        taste a little, so that it may look like a visit to my house!"}
}
{
}

\Verse{40}
{
    \zA{说着翻身起来，将两只手呵了两口，便伸向黛玉膈 肢窝内两胁下乱挠。 黛玉素性触痒不禁，见宝玉两手伸来乱挠，便笑的喘不过气来。
        口里说：“宝玉!你再闹，我就恼了！” 宝玉方住了手，笑问道：“你还说这些不说了？” 黛玉笑道：“再不敢了。”
        一面理鬓笑道：“我有奇香，你有‘暖香’没有？” 宝玉见问，一时解不来，因问： “什么‘暖香’？”
        黛玉点头笑叹道：“蠢才，蠢才!你有玉，人家就有金来配你； 人家有‘冷香’，你就没有‘暖香’去配他？”}
}
{
    \eA{As she said this, she forthwith took several seeds of the fir-cone, and
        cracking off the thin skin, she placed them in a handkerchief and presented
        them to Pao-yü. But Pao-yü, espying that Hsi Jen's two eyes were slightly
        red, and that the powder was shiny and moist, quietly therefore inquired of
        Hsi Jen, "Why do you cry for no rhyme or reason?" "Why should I cry?" Hsi
        Jen laughed;}
}
{
}

\Verse{41}
{
    \zA{宝玉方听出来，因笑道：“方才告 饶，如今更说狠了！” 说着又要伸手。 黛玉忙笑道：“好哥哥，我可不敢了。” 宝 玉笑道：“饶你不难，只把袖子我闻一闻。”
        说着便拉了袖子笼在面上，闻个不住。 黛玉夺了手道：“这可该去了。” 宝玉笑道：“要去不能。 咱们斯斯文文的躺着说 话儿。”
        说着复又躺下，黛玉也躺下，用绢子盖上脸。 宝玉有一搭没一搭的说些鬼话，黛玉总不理。 宝玉问他几岁上京，路上见何景
        致，扬州有何古迹，土俗民风如何，黛玉不答。}
}
{
    \eA{"something just got into my eyes and I rubbed them." By these means she
        readily managed to evade detection; but seeing that Pao-yü wore a deep red
        archery-sleeved pelisse, ornamented with gold dragons, and lined with fur
        from foxes' ribs and a grey sable fur surtout with a fringe round the
        border. "What! have you," she asked, "put on again your new clothes for?
        specially to come here?}
}
{
}

\Verse{42}
{
    \zA{宝玉只怕他睡出病来，便哄他道： “嗳哟!你们扬州衙门里有一件大故事，你可知道么？” 黛玉见他说的郑重，又且 正言厉色，只当是真事，因问：“什么事？”
        宝玉见问，便忍着笑顺口诌道：“扬 州有一座黛山，山上有个林子洞。” 黛玉笑道：“这就扯谎，自来也没听见这山。” 宝玉道：“天下山水多着呢，你那里都知道？
        等我说完了你再批评。” 黛玉道：“你 说。” 宝玉又诌道：“林子洞里原来有一群耗子精。}
}
{
    \eA{and didn't they inquire of you where you were going?" "I had changed,"
        Pao-yü explained with a grin, "as Mr. Chen had invited me to go over and
        look at the play." "Well, sit a while and then go back;" Hsi Jen continued
        as she nodded her head; "for this isn't the place for you to come to!"
        "You'd better be going home now," Pao-yü suggested smirkingly; "where I've
        again kept something good for you."}
}
{
}

\Verse{43}
{
    \zA{那一年腊月初七老耗子升座议事， 说：‘明儿是腊八儿了，世上的人都熬腊八粥，如今我们洞里果品短少，须得趁此 打劫些个来才好。’
        乃拔令箭一枝，遣了个能干小耗子去打听。 小耗子回报：‘各 处都打听了，惟有山下庙里果米最多。’ 老耗子便问：‘米有几样？ 果有几品？’
        小耗子道：‘米豆成仓。 果品却只有五样：一是红枣，二是栗子，三是落花生，四 是菱角，五是香芋。’ 老耗子听了大喜，即时拔了一枝令箭，问：‘谁去偷米？’
        一个耗子便接令去偷米。}
}
{
    \eA{"Gently," smiled Hsi Jen, "for were you to let them hear, what figure would
        we cut?" And with these, words, she put out her hand and unclasping from
        Pao-yü's neck the jade of Spiritual Perception, she faced her cousins and
        remarked exultingly. "Here! see for yourselves; look at this and learn! When
        I repeatedly talked about it, you all thought it extraordinary, and were
        anxious to have a glance at it;}
}
{
}

\Verse{44}
{
    \zA{又拔令箭问：‘谁去偷豆？’ 又一个耗子接令去偷豆。 然 后一一的都各领令去了。 只剩下香芋。 因又拔令箭问：‘谁去偷香芋？’ 只见一个
        极小极弱的小耗子应道：‘我愿去偷香芋。” 老耗子和众耗见他这样，恐他不谙练， 又怯懦无力，不准他去。 小耗子道：‘我虽年小身弱，却是法术无边，口齿伶俐，
        机谋深远。 这一去，管比他们偷的还巧呢！’ 众耗子忙问：‘怎么比他们巧呢？’}
}
{
    \eA{to-day, you may gaze on it with all your might, for whatever precious thing
        you may by and by come to see will really never excel such an object as
        this!"}
}
{
}

\Verse{45}
{
    \zA{小耗子道：‘我不学他们直偷，我只摇身一变，也变成个香芋，滚在香芋堆里，叫 人瞧不出来，却暗暗儿的搬运，渐渐的就搬运尽了：这不比直偷硬取的巧吗？’ 众
        耗子听了，都说：‘妙却妙，只是不知怎么变？ 你先变个我们瞧瞧。’ 小耗子听了， 笑道：‘这个不难，等我变来。’ 说毕，摇身说：‘变。’
        竟变了一个最标致美貌 的一位小姐。 众耗子忙笑说：‘错了，错了!原说变果子，怎么变出个小姐来了呢？’}
}
{
    \eA{When she had finished speaking, she handed it over to them, and after they
        had passed it round for inspection, she again fastened it properly on
        Pao-yü's neck, and also bade her brother go and hire a small carriage, or
        engage a small chair, and escort Pao-yü back home. "If I see him back," Hua
        Tzu-fang remarked, "there would be no harm, were he even to ride his horse!"
        "It isn't because of harm," Hsi Jen replied;}
}
{
}

\Verse{46}
{
    \zA{小耗子现了形笑道：‘我说你们没见世面，只认得这果子是香芋，却不知盐课林老 爷的小姐才是真正的“香玉”呢！’”
        黛玉听了，翻身爬起来，按着宝玉笑道：“我把你这个烂了嘴的!我就知道你 是编派我呢。” 说着便拧。 宝玉连连央告：“好妹妹，饶了我罢，再不敢了。 我因
        为闻见你的香气，忽然想起这个故典来。” 黛玉笑道：“饶骂了人，你还说是故典 呢。” 一语未了，只见宝钗走来，笑问：“谁说故典呢？ 我也听听。”}
}
{
    \eA{"but because he may come across some one from the house." Hua Tzu-fang
        promptly went and bespoke a small chair; and when it came to the door, the
        whole party could not very well detain him, and they of course had to see
        Pao-yü out of the house; while Hsi Jen, on the other hand, snatched a few
        fruits and gave them to Ming Yen;}
}
{
}

\Verse{47}
{
    \zA{黛玉忙让坐， 笑道：“你瞧瞧，还有谁？ 他饶骂了，还说是故典。” 宝钗笑道：“哦!是宝兄弟哟! 怪不得他。
        他肚子里的故典本来多么!就只是可惜一件，该用故典的时候儿他就偏 忘了。 有今儿记得的，前儿夜里的芭蕉诗就该记得呀，眼面前儿的倒想不起来。 别
        人冷的了不得，他只是出汗。 这会子偏又有了记性了！” 黛玉听了笑道：“阿弥陀 佛!到底是我的好姐姐。 你一般也遇见对子了。 可知一还一报，不爽不错的。”}
}
{
    \eA{and as she at the same time pressed in his hand several cash to buy crackers
        with to let off, she enjoined him not to tell any one as he himself would
        likewise incur blame. As she uttered these words, she straightway escorted
        Pao-yü as far as outside the door, from whence having seen him mount into
        the sedan chair, she dropped the curtain;}
}
{
}

\Verse{48}
{
    \zA{刚 说到这里，只听宝玉房中一片声吵嚷起来。 未知何事，下回分解。 (本章完)}
}
{
    \eA{whereupon Ming Yen and her brother, the two of them, led the horses and
        followed behind in his wake. Upon reaching the street where the Ning mansion
        was situated, Ming Yen told the chair to halt, and said to Hua Tzu-fang,
        "It's advisable that I should again go, with Mr. Secundus, into the Eastern
        mansion, to show ourselves before we can safely betake ourselves home; for
        if we don't, people will suspect!"}
}
{
}

\Verse{49}
{
    \zA{}
}
{
    \eA{Hua Tzu-fang, upon hearing that there was good reason in what he said,
        promptly clasped Pao-yü out of the chair and put him on the horse, whereupon
        after Pao-yü smilingly remarked: "Excuse me for the trouble I've surely put
        you to," they forthwith entered again by the back gate; but putting aside
        all details, we will now confine ourselves to Pao-yü. After he had walked
        out of the door, the several waiting-maids in his apartments played and
        laughed with greater zest and with less restraint. Some there were who
        played at chess, others who threw the dice or had a game of cards;}
}
{
}

\Verse{50}
{
    \zA{}
}
{
    \eA{and they covered the whole floor with the shells of melon-seeds they were
        cracking, when dame Li, his nurse, happened to come in, propping herself on
        a staff, to pay her respects and to see Pao-yü, and perceiving that Pao-yü
        was not at home and that the servant-girls were only bent upon romping, she
        felt intensely disgusted. "Since I've left this place," she therefore
        exclaimed with a sigh, "and don't often come here, you've become more and
        more unmannerly; while the other nurse does still less than ever venture to
        expostulate with you; Pao-yü is like a candlestick eighty feet high,
        shedding light on others, and throwing none upon himself! All he knows is to
        look down upon people as being filthy;}
}
{
}

\Verse{51}
{
    \zA{}
}
{
    \eA{and yet this is his room and he allows you to put it topsy-turvey, and to
        become more and more unmindful of decorum!" These servant-girls were well
        aware that Pao-yü was not particular in these respects, and that in the next
        place nurse Li, having pleaded old age, resigned her place and gone home,
        had nowadays no control over them, so that they simply gave their minds to
        romping and joking, and paid no heed whatever to her. Nurse Li however still
        kept on asking about Pao-yü, "How much rice he now ate at one meal?}
}
{
}

\Verse{52}
{
    \zA{}
}
{
    \eA{and at what time he went to sleep?" to which questions, the servant-girls
        replied quite at random; some there being too who observed: "What a dreadful
        despicable old thing she is!" "In this covered bowl," she continued to
        inquire, "is cream, and why not give it to me to eat?" and having concluded
        these words, she took it up and there and then began eating it. "Be quick,
        and leave it alone!" a servant-girl expostulated, "that, he said, was kept
        in order to be given to Hsi Jen; and on his return, when he again gets into
        a huff, you, old lady, must, on your own motion, confess to having eaten it,
        and not involve us in any way as to have to bear his resentment." Nurse Li,
        at these words, felt both angry and ashamed. "I can't believe," she
        forthwith remarked, "that he has become so bad at heart!}
}
{
}

\Verse{53}
{
    \zA{}
}
{
    \eA{Not to speak of the milk I've had, I have, in fact every right to even
        something more expensive than this; for is it likely that he holds Hsi Jen
        dearer than myself? It can't forsooth be that he doesn't bear in mind how
        that I've brought him up to be a big man, and how that he has eaten my blood
        transformed into milk and grown up to this age! and will be because I'm now
        having a bowl of milk of his be angry on that score! I shall, yes, eat it,
        and we'll see what he'll do! I don't know what you people think of Hsi Jen,
        but she was a lowbred girl, whom I've with my own hands raised up! and what
        fine object indeed was she!" As she spoke, she flew into a temper, and
        taking the cream she drank the whole of it. "They don't know how to speak
        properly!"}
}
{
}

\Verse{54}
{
    \zA{}
}
{
    \eA{another servant-girl interposed sarcastically, "and it's no wonder that you,
        old lady, should get angry! Pao-yü still sends you, venerable dame, presents
        as a proof of his gratitude, and is it possible that he will feel displeased
        for such a thing like this?" "You girls shouldn't also pretend to be artful
        flatterers to cajole me!" nurse Li added; "do you imagine that I'm not aware
        of the dismissal, the other day, of Hsi Hsüeh, on account of a cup of tea?
        and as it's clear enough that I've incurred blame, I'll come by and by and
        receive it!" Having said this, she went off in a dudgeon, but not a long
        interval elapsed before Pao-yü returned, and gave orders to go and fetch Hsi
        Jen;}
}
{
}

\Verse{55}
{
    \zA{}
}
{
    \eA{and perceiving Ching Ling reclining on the bed perfectly still: "I presume
        she's ill," Pao-yü felt constrained to inquire, "or if she isn't ill, she
        must have lost at cards." "Not so!" observed Chiu Wen; "she had been a
        winner, but dame Li came in quite casually and muddled her so that she lost;
        and angry at this she rushed off to sleep." "Don't place yourselves," Pao-yü
        smiled, "on the same footing as nurse Li, and if you were to let her alone,
        everything will be all right." These words were still on his lips when Hsi
        Jen arrived. After the mutual salutations, Hsi Jen went on to ask of Pao-yü:
        "Where did you have your repast? and what time did you come back?"}
}
{
}

\Verse{56}
{
    \zA{}
}
{
    \eA{and to present likewise, on behalf of her mother and sister, her compliments
        to all the girls, who were her companions. In a short while, she changed her
        costume and divested herself of her fineries, and Pao-yü bade them fetch the
        cream. "Nurse Li has eaten it," the servant-girls rejoined, and as Pao-yü
        was on the point of making some remark Hsi Jen hastened to interfere,
        laughing the while; "Is it really this that you had kept for me? many thanks
        for the trouble; the other day, when I had some, I found it very toothsome,
        but after I had partaken of it, I got a pain in the stomach, and was so much
        upset, that it was only after I had brought it all up that I felt all right.
        So it's as well that she has had it, for, had it been kept here, it would
        have been wasted all for no use!}
}
{
}

\Verse{57}
{
    \zA{}
}
{
    \eA{What I fancy are dry chestnuts; and while you clean a few for me, I'll go
        and lay the bed!" Pao-yü upon hearing these words credited them as true, so
        that he discarded all thought of the cream and fetched the chestnuts, which
        he, with his own hands, selected and pealed. Perceiving at the same time
        that none of the party were present in the room, he put on a smile and
        inquired of Hsi Jen: "Who were those persons dressed in red to day?"
        "They're my two cousins on my mother's side," Hsi Jen explained, and hearing
        this, Pao-yü sang their praise as he heaved a couple of sighs. "What are you
        sighing for?" Hsi Jen remarked. "I know the secret reasons of your heart;}
}
{
}

\Verse{58}
{
    \zA{}
}
{
    \eA{it's I fancy because she isn't fit to wear red!" "It isn't that," Pao-yü
        protested smilingly, "it isn't that; if such a person as that isn't good
        enough to be dressed in red, who would forsooth presume to wear it? It's
        because I find her so really lovely! and if we could, after all, manage to
        get her into our family, how nice it would be then!" Hsi Jen gave a sardonic
        smile. "That it's my own fate to be a slave doesn't matter, but is it likely
        that the destiny of even my very relatives could be to become one and all of
        them bond servants? But you should certainly set your choice upon some
        really beautiful girl, for she would in that case be good enough to enter
        your house." "Here you are again with your touchiness!"}
}
{
}

\Verse{59}
{
    \zA{}
}
{
    \eA{Pao-yü eagerly exclaimed smiling, "if I said that she should come to our
        house, does it necessarily imply that she should be a servant? and wouldn't
        it do were I to mention that she should come as a relative!" "That too
        couldn't exalt her to be a fit match for you!" rejoined Hsi Jen; but Pao-yü
        being loth to continue the conversation, simply busied himself with cleaning
        the chestnuts. "How is it you utter not a word?" Hsi Jen laughed; "I expect
        it's because I just offended you by my inconsiderate talk! But if by and by
        you have your purpose fixed on it, just spend a few ounces of silver to
        purchase them with, and bring them in and have done!" "How would you have
        one make any reply?" Pao-yü smilingly rejoined; "all I did was to extol her
        charms;}
}
{
}

\Verse{60}
{
    \zA{}
}
{
    \eA{for she's really fit to have been born in a deep hall and spacious court as
        this; and it isn't for such foul things as myself and others to contrariwise
        spend our days in this place!" "Though deprived of this good fortune," Hsi
        Jen explained, "she's nevertheless also petted and indulged and the jewel of
        my maternal uncle and my aunt! She's now seventeen years of age, and
        everything in the way of trousseau has been got ready, and she's to get
        married next year." Upon hearing the two words "get married," he could not
        repress himself from again ejaculating: "Hai hai!"}
}
{
}

\Verse{61}
{
    \zA{}
}
{
    \eA{but while he was in an unhappy frame of mind, he once more heard Hsi Jen
        remark as she heaved a sigh: "Ever since I've come here, we cousins haven't
        all these years been able to get to live together, and now that I'm about to
        return home, they, on the other hand, will all be gone!" Pao-yü, realising
        that there lurked in this remark some meaning or other, was suddenly so
        taken aback that dropping the chestnuts, he inquired: "How is it that you
        now want to go back?" "I was present to-day," Hsi Jen explained, "when
        mother and brother held consultation together, and they bade me be patient
        for another year, and that next year they'll come up and redeem me out of
        service!"}
}
{
}

\Verse{62}
{
    \zA{}
}
{
    \eA{Pao-yü, at these words, felt the more distressed. "Why do they want to
        redeem you?" he consequently asked. "This is a strange question!" Hsi Jen
        retorted, "for I can't really be treated as if I were the issue born in this
        homestead of yours! All the members of my family are elsewhere, and there's
        only myself in this place, so that how could I end my days here?" "If I
        don't let you go, it will verily be difficult for you to get away!" Pao-yü
        replied. "There has never been such a principle of action!" urged Hsi Jen;
        "even in the imperial palace itself, there's a fixed rule, by which possibly
        every certain number of years a selection (of those who have to go takes
        place), and every certain number of years a new batch enters;}
}
{
}

\Verse{63}
{
    \zA{}
}
{
    \eA{and there's no such practice as that of keeping people for ever; not to
        speak of your own home." Pao-yü realised, after reflection, that she, in
        point of fact, was right, and he went on to observe: "Should the old lady
        not give you your release, it will be impossible for you to get off." "Why
        shouldn't she release me?" Hsi Jen questioned. "Am I really so very
        extraordinary a person as to have perchance made such an impression upon her
        venerable ladyship and my lady that they will be positive in not letting me
        go? They may, in all likelihood, give my family some more ounces of silver
        to keep me here; that possibly may come about. But, in truth, I'm also a
        person of the most ordinary run, and there are many more superior to me, yea
        very many!}
}
{
}

\Verse{64}
{
    \zA{}
}
{
    \eA{Ever since my youth up, I've been in her old ladyship's service; first by
        waiting upon Miss Shih for several years, and recently by being in
        attendance upon you for another term of years; and now that our people will
        come to redeem me, I should, as a matter of right, be told to go. My idea is
        that even the very redemption money won't be accepted, and that they will
        display such grace as to let me go at once. And, as for being told that I
        can't be allowed to go as I'm so diligent in my service to you, that's a
        thing that can on no account come about! My faithful attendance is an
        obligation of my duties, and is no exceptional service! and when I'm gone
        you'll again have some other faithful attendant, and it isn't likely that
        when I'm no more here, you'll find it impracticable to obtain one!"}
}
{
}

\Verse{65}
{
    \zA{}
}
{
    \eA{After Pao-yü had listened to these various arguments, which proved the
        reasonableness of her going and the unreasonableness of any detention, he
        felt his heart more than ever a prey to distress. "In spite of all you say,"
        he therefore continued, "the sole desire of my heart is to detain you; and I
        have no doubt but that the old lady will speak to your mother about it; and
        if she were to give your mother ample money, she'll, of course, not feel as
        if she could very well with any decency take you home!" "My mother won't
        naturally have the audacity to be headstrong!" Hsi Jen ventured, "not to
        speak besides of the nice things, which may be told her and the lots of
        money she may, in addition, be given;}
}
{
}

\Verse{66}
{
    \zA{}
}
{
    \eA{but were she even not to be paid any compliments, and not so much as a
        single cash given her, she won't, if you set your mind upon keeping me here,
        presume not to comply with your wishes, were it also against my inclination.
        One thing however; our family would never rely upon prestige, and trust upon
        honorability to do anything so domineering as this! for this isn't like
        anything else, which, because you take a fancy to it, a hundred per cent
        profit can be added, and it obtained for you! This action can be well taken
        if the seller doesn't suffer loss! But in the present instance, were they to
        keep me back for no rhyme or reason, it would also be of no benefit to
        yourself; on the contrary, they would be instrumental in keeping us blood
        relatives far apart;}
}
{
}

\Verse{67}
{
    \zA{}
}
{
    \eA{a thing the like of which, I feel positive that dowager lady Chia and my
        lady will never do!" After lending an ear to this argument, Pao-yü cogitated
        within himself for a while. "From what you say," he then observed, "when you
        say you'll go, it means that you'll go for certain!" "Yes, that I'll go for
        certain," Hsi Jen rejoined. "Who would have anticipated," Pao-yü, after
        these words, mused in his own heart, "that a person like her would have
        shown such little sense of gratitude, and such a lack of respect! Had I," he
        then remarked aloud with a sigh, "been aware, at an early date, that your
        whole wish would have been to go, I wouldn't, in that case, have brought you
        over! But when you're away, I shall remain alone, a solitary spirit!"}
}
{
}

\Verse{68}
{
    \zA{}
}
{
    \eA{As he spoke, he lost control over his temper, and, getting into bed, he went
        to sleep. The fact is that when Hsi Jen had been at home, and she heard her
        mother and brother express their intention of redeeming her back, she there
        and then observed that were she even at the point of death, she would not
        return home. "When in past days," she had argued, "you had no rice to eat,
        there remained myself, who was still worth several taels; and hadn't I urged
        you to sell me, wouldn't I have seen both father and mother die of
        starvation under my very eyes? and you've now had the good fortune of
        selling me into this place, where I'm fed and clothed just like a mistress,
        and where I'm not beaten by day, nor abused by night!}
}
{
}

\Verse{69}
{
    \zA{}
}
{
    \eA{Besides, though now father be no more, you two have anyhow by putting things
        straight again, so adjusted the family estate that it has resumed its
        primitive condition. And were you, in fact, still in straitened
        circumstances, and you could by redeeming me back, make again some more
        money, that would be well and good; but the truth is that there's no such
        need, and what would be the use for you to redeem me at such a time as this?
        You should temporarily treat me as dead and gone, and shouldn't again recall
        any idea of redeeming me!"}
}
{
}

\Verse{70}
{
    \zA{}
}
{
    \eA{Having in consequence indulged in a loud fit of crying, her mother and
        brother resolved, when they perceived her in this determined frame of mind,
        that for a fact there was no need for her to come out of service. What is
        more they had sold her under contract until death, in the distinct reliance
        that the Chia family, charitable and generous a family as it was, would,
        possibly, after no more than a few entreaties, make them a present of her
        person as well as the purchase money. In the second place, never had they in
        the Chia mansion ill-used any of those below; there being always plenty of
        grace and little of imperiousness.}
}
{
}

\Verse{71}
{
    \zA{}
}
{
    \eA{Besides, the servant-girls, who acted as personal attendants in the
        apartments of the old as well as of the young, were treated so far unlike
        the whole body of domestics in the household that the daughters even of an
        ordinary and penniless parentage could not have been so looked up to. And
        these considerations induced both the mother as well as her son to at once
        dispel the intention and not to redeem her, and when Pao-yü had subsequently
        paid them an unexpected visit, and the two of them (Pao-yü and Hsi Jen) were
        seen to be also on such terms, the mother and her son obtained a clearer
        insight into their relations, and still one more burden (which had pressed
        on their mind) fell to the ground, and as besides this was a contingency,
        which they had never reckoned upon, they both composed their hearts, and did
        not again entertain any idea of ransoming her.}
}
{
}

\Verse{72}
{
    \zA{}
}
{
    \eA{It must be noticed moreover that Hsi Jen had ever since her youth not been
        blind to the fact that Pao-yü had an extraordinary temperament, that he was
        self-willed and perverse, far even in excess of all young lads, and that he
        had, in addition, a good many peculiarities and many unspeakable defects.
        And as of late he had placed such reliance in the fond love of his
        grandmother that his father and mother even could not exercise any extreme
        control over him, he had become so much the more remiss, dissolute, selfish
        and unconcerned, not taking the least pleasure in what was proper, that she
        felt convinced, whenever she entertained the idea of tendering him advice,
        that he would not listen to her.}
}
{
}

\Verse{73}
{
    \zA{}
}
{
    \eA{On this day, by a strange coincidence, came about the discussion respecting
        her ransom, and she designedly made use, in the first instance, of deception
        with a view to ascertain his feelings, to suppress his temper, and to be
        able subsequently to extend to him some words of admonition; and when she
        perceived that Pao-yü had now silently gone to sleep, she knew that his
        feelings could not brook the idea of her return and that his temper had
        already subsided.}
}
{
}

\Verse{74}
{
    \zA{}
}
{
    \eA{She had never had, as far as she was concerned, any desire of eating
        chestnuts, but as she feared lest, on account of the cream, some trouble
        might arise, which might again lead to the same results as when Hsi Hsüeh
        drank the tea, she consequently made use of the pretence that she fancied
        chestnuts, in order to put off Pao-yü from alluding (to the cream) and to
        bring the matter speedily to an end. But telling forthwith the young
        waiting-maids to take the chestnuts away and eat them, she herself came and
        pushed Pao-yü; but at the sight of Pao-yü with the traces of tears on his
        face, she at once put on a smiling expression and said: "What's there in
        this to wound your heart? If you positively do wish to keep me, I shall, of
        course, not go away!"}
}
{
}

\Verse{75}
{
    \zA{}
}
{
    \eA{Pao-yü noticed that these words contained some hidden purpose, and readily
        observed: "Do go on and tell me what else I can do to succeed in keeping you
        here, for of my own self I find it indeed difficult to say how!" "Of our
        friendliness all along," Hsi Jen smilingly rejoined, "there's naturally no
        need to speak; but, if you have this day made up your mind to retain me
        here, it isn't through this friendship that you'll succeed in doing so. But
        I'll go on and mention three distinct conditions, and, if you really do
        accede to my wishes, you'll then have shown an earnest desire to keep me
        here, and I won't go, were even a sword to be laid on my neck!"}
}
{
}

\Verse{76}
{
    \zA{}
}
{
    \eA{"Do tell me what these conditions are," Pao-yü pressed her with alacrity, as
        he smiled, "and I'll assent to one and all. My dear sister, my own dear
        sister, not to speak of two or three, but even two or three hundred of them
        I'm quite ready to accept. All I entreat you is that you and all of you
        should combine to watch over me and take care of me, until some day when I
        shall be transformed into flying ashes; but flying ashes are, after all, not
        opportune, as they have form and substance and they likewise possess sense,
        but until I've been metamorphosed into a streak of subtle smoke. And when
        the wind shall have with one puff dispelled me, all of you then will be
        unable to attend to me, just as much as I myself won't be able to heed you.}
}
{
}

\Verse{77}
{
    \zA{}
}
{
    \eA{You will, when that time comes, let me go where I please, as I'll let you
        speed where you choose to go!" These words so harassed Hsi Jen that she
        hastened to put her hand over his mouth. "Speak decently," she said; "I was
        on account of this just about to admonish you, and now here you are uttering
        all this still more loathsome trash." "I won't utter these words again,"
        Pao-yü eagerly added. "This is the first fault that you must change," Hsi
        Jen replied. "I'll amend," Pao-yü observed, "and if I say anything of the
        kind again you can wring my mouth; but what else is there?" "The second
        thing is this," Hsi Jen explained; "whether you really like to study or
        whether you only pretend to like study is immaterial;}
}
{
}

\Verse{78}
{
    \zA{}
}
{
    \eA{but you should, when you are in the presence of master, or in the presence
        of any one else, not do nothing else than find fault with people and make
        fun of them, but behave just as if you were genuinely fond of study, so that
        you shouldn't besides provoke your father so much to anger, and that he
        should before others have also a chance of saying something! 'In my family,'
        he reflects within himself, 'generation after generation has been fond of
        books, but ever since I've had you, you haven't accomplished my
        expectations, and not only is it that you don't care about reading
        books,'--and this has already filled his heart with anger and
        vexation,--'but both before my face and behind my back, you utter all that
        stuff and nonsense, and give those persons, who have, through their
        knowledge of letters, attained high offices, the nickname of the "the
        salaried worms."}
}
{
}

\Verse{79}
{
    \zA{}
}
{
    \eA{You also uphold that there's no work exclusive (of the book where appears)
        "fathom spotless virtue;" and that all other books consist of foolish
        compilations, which owe their origin to former authors, who, unable
        themselves to expound the writings of Confucius, readily struck a new line
        and invented original notions.' Now with words like these, how can one
        wonder if master loses all patience, and if he does from time to time give
        you a thrashing! and what do you make other people think of you?" "I won't
        say these things again," Pao-yü laughingly protested, "these are the
        reckless and silly absurdities of a time when I was young and had no idea of
        the height of the heavens and the thickness of the earth;}
}
{
}

\Verse{80}
{
    \zA{}
}
{
    \eA{but I'll now no more repeat them. What else is there besides?" "It isn't
        right that you should sneer at the bonzes and vilify the Taoist priests, nor
        mix cosmetics or prepare rouge," Hsi Jen continued; "but there's still
        another thing more important, you shouldn't again indulge the bad habits of
        licking the cosmetic, applied by people on their lips, nor be fond of (girls
        dressed) in red!" "I'll change in all this," Pao-yü added by way of
        rejoinder; "I'll change in all this; and if there's anything more be quick
        and tell me." "There's nothing more," Hsi Jen observed; "but you must in
        everything exercise a little more diligence, and not indulge your caprices
        and allow your wishes to run riot, and you'll be all right.}
}
{
}

\Verse{81}
{
    \zA{}
}
{
    \eA{And should you comply to all these things in real earnest, you couldn't
        carry me out, even in a chair with eight bearers." "Well, if you do stay in
        here long enough," Pao-yü remarked with a smile, "there's no fear as to your
        not having an eight-bearer-chair to sit in!" Hsi Jen gave a sardonic grin.
        "I don't care much about it," she replied; "and were I even to have such
        good fortune, I couldn't enjoy such a right. But allowing I could sit in
        one, there would be no pleasure in it!" While these two were chatting, they
        saw Ch'iu Wen walk in. "It's the third watch of the night," she observed,
        "and you should go to sleep. Just a few moments back your grandmother lady
        Chia and our lady sent a nurse to ask about you, and I replied that you were
        asleep."}
}
{
}

\Verse{82}
{
    \zA{}
}
{
    \eA{Pao-yü bade her fetch a watch, and upon looking at the time, he found indeed
        that the hand was pointing at ten; whereupon rinsing his mouth again and
        loosening his clothes, he retired to rest, where we will leave him without
        any further comment. The next day, Hsi Jen got up as soon as it was dawn,
        feeling her body heavy, her head sore, her eyes swollen, and her limbs
        burning like fire. She managed however at first to keep up, an effort though
        it was, but as subsequently she was unable to endure the strain, and all she
        felt disposed to do was to recline, she therefore lay down in her clothes on
        the stove-couch.}
}
{
}

\Verse{83}
{
    \zA{}
}
{
    \eA{Pao-yü hastened to tell dowager lady Chia, and the doctor was sent for, who,
        upon feeling her pulse and diagnosing her complaint, declared that there was
        nothing else the matter with her than a chill, which she had suddenly
        contracted, that after she had taken a dose or two of medicine, it would be
        dispelled, and that she would be quite well. After he had written the
        prescription and taken his departure, some one was despatched to fetch the
        medicines, which when brought were properly decocted. As soon as she had
        swallowed a dose, Pao-yü bade her cover herself with her bed-clothes so as
        to bring on perspiration; while he himself came into Tai-yü's room to look
        her up.}
}
{
}

\Verse{84}
{
    \zA{}
}
{
    \eA{Tai-yü was at this time quite alone, reclining on her bed having a midday
        siesta, and the waiting-maids having all gone out to attend to whatever they
        pleased, the whole room was plunged in stillness and silence. Pao-yü raised
        the embroidered soft thread portiere and walked in; and upon espying Tai-yü
        in the room fast asleep, he hurriedly approached her and pushing her: "Dear
        cousin," he said, "you've just had your meal, and are you asleep already?"
        and he kept on calling "Tai-yü" till he woke her out of her sleep.}
}
{
}

\Verse{85}
{
    \zA{}
}
{
    \eA{Perceiving that it was Pao-yü, "You had better go for a stroll," Tai-yü
        urged, "for the day before yesterday I was disturbed the whole night, and up
        to this day I haven't had rest enough to get over the fatigue. My whole body
        feels languid and sore." "This languor and soreness," Pao-yü rejoined, "are
        of no consequence; but if you go on sleeping you'll be feeling very ill; so
        I'll try and distract you, and when we've dispelled this lassitude, you'll
        be all right." Tai-yü closed her eyes. "I don't feel any lassitude," she
        explained, "all I want is a little rest; and you had better go elsewhere and
        come back after romping about for a while." "Where can I go?" Pao-yü asked
        as he pushed her. "I'm quite sick and tired of seeing the others."}
}
{
}

\Verse{86}
{
    \zA{}
}
{
    \eA{At these words, Tai-yü burst out laughing with a sound of Ch'ih. "Well!
        since you wish to remain here," she added, "go over there and sit down
        quietly, and let's have a chat." "I'll also recline," Pao-yü suggested.
        "Well, then, recline!" Tai-yü assented. "There's no pillow," observed
        Pao-yü, "so let us lie on the same pillow." "What nonsense!" Tai-yü urged,
        "aren't those pillows outside? get one and lie on it." Pao-yü walked into
        the outer apartment, and having looked about him, he returned and remarked
        with a smile: "I don't want those, they may be, for aught I know, some dirty
        old hag's." Tai-yü at this remark opened her eyes wide, and as she raised
        herself up: "You're really," she exclaimed laughingly, "the evil star of my
        existence! here, please recline on this pillow!"}
}
{
}

\Verse{87}
{
    \zA{}
}
{
    \eA{and as she uttered these words, she pushed her own pillow towards Pao-yü,
        and, getting up she went and fetched another of her own, upon which she lay
        her head in such a way that both of them then reclined opposite to each
        other. But Tai-yü, upon turning up her eyes and looking, espied on Pao-yü's
        cheek on the left side of his face, a spot of blood about the size of a
        button, and speedily bending her body, she drew near to him, and rubbing it
        with her hand, she scrutinised it closely. "Whose nail," she went on to
        inquire, "has scratched this open?"}
}
{
}

\Verse{88}
{
    \zA{}
}
{
    \eA{Pao-yü with his body still reclining withdrew from her reach, and as he did
        so, he answered with a smile: "It isn't a scratch; it must, I presume, be
        simply a drop, which bespattered my cheek when I was just now mixing and
        clarifying the cosmetic paste for them." Saying this, he tried to get at his
        handkerchief to wipe it off; but Tai-yü used her own and rubbed it clean for
        him, while she observed: "Do you still give your mind to such things? attend
        to them you may; but must you carry about you a placard (to make it public)?}
}
{
}

\Verse{89}
{
    \zA{}
}
{
    \eA{Though uncle mayn't see it, were others to notice it, they would treat it as
        a strange occurrence and a novel bit of news, and go and tell him to curry
        favour, and when it has reached uncle's ear, we shall all again not come out
        clean, and provoke him to anger." Pao-yü did not in the least heed what she
        said, being intent upon smelling a subtle scent which, in point of fact,
        emanated from Tai-yü's sleeve, and when inhaled inebriated the soul and
        paralysed the bones. With a snatch, Pao-yü laid hold of Tai-yü's sleeve
        meaning to see what object was concealed in it; but Tai-yü smilingly
        expostulated: "At such a time as this," she said, "who keeps scents about
        one?"}
}
{
}

\Verse{90}
{
    \zA{}
}
{
    \eA{"Well, in that case," Pao-yü rejoined with a smirking face, "where does this
        scent come from?" "I myself don't know," Tai-yü replied; "I presume it must
        be, there's no saying, some scent in the press which has impregnated the
        clothes." "It doesn't follow," Pao-yü added, as he shook his head; "the
        fumes of this smell are very peculiar, and don't resemble the perfume of
        scent-bottles, scent-balls, or scented satchets!" "Is it likely that I have,
        like others, Buddhistic disciples," Tai-yü asked laughing ironically, "or
        worthies to give me novel kinds of scents? But supposing there is about me
        some peculiar scent, I haven't, at all events, any older or younger brothers
        to get the flowers, buds, dew, and snow, and concoct any for me; all I have
        are those common scents, that's all."}
}
{
}

\Verse{91}
{
    \zA{}
}
{
    \eA{"Whenever I utter any single remark," Pao-yü urged with a grin, "you at once
        bring up all these insinuations; but unless I deal with you severely, you'll
        never know what stuff I'm made of; but from henceforth I'll no more show you
        any grace!" As he spoke, he turned himself over, and raising himself, he
        puffed a couple of breaths into both his hands, and hastily stretching them
        out, he tickled Tai-yü promiscuously under her armpits, and along both
        sides. Tai-yü had never been able to stand tickling, so that when Pao-yü put
        out his two hands and tickled her violently, she forthwith giggled to such
        an extent that she could scarcely gasp for breath.}
}
{
}

\Verse{92}
{
    \zA{}
}
{
    \eA{"If you still go on teasing me," she shouted, "I'll get angry with you!"
        Pao-yü then kept his hands off, and as he laughed, "Tell me," he asked,
        "will you again come out with all those words or not?" "I daren't do it
        again," Tai-yü smiled and adjusted her hair; adding with another laugh: "I
        may have peculiar scents, but have you any 'warm' scents?" Pao-yü at this
        question, could not for a time unfold its meaning: "What 'warm' scent?" he
        therefore asked. Tai-yü nodded her head and smiled deridingly. "How stupid!
        what a fool!" she sighed; "you have jade, and another person has gold to
        match with you, and if some one has 'cold' scent, haven't you any 'warm'
        scent as a set-off?" Pao-yü at this stage alone understood the import of her
        remark.}
}
{
}

\Verse{93}
{
    \zA{}
}
{
    \eA{"A short while back you craved for mercy," Pao-yü observed smilingly, "and
        here you are now going on talking worse than ever;" and as he spoke he again
        put out his hands. "Dear cousin," Tai-yü speedily implored with a smirk, "I
        won't venture to do it again." "As for letting you off," Pao-yü remarked
        laughing, "I'll readily let you off, but do allow me to take your sleeve and
        smell it!" and while uttering these words, he hastily pulled the sleeve, and
        pressing it against his face, kept on smelling it incessantly, whereupon
        Tai-yü drew her hand away and urged: "You must be going now!"}
}
{
}

\Verse{94}
{
    \zA{}
}
{
    \eA{"Though you may wish me to go, I can't," Pao-yü smiled, "so let us now lie
        down with all propriety and have a chat," laying himself down again, as he
        spoke, while Tai-yü likewise reclined, and covered her face with her
        handkerchief. Pao-yü in a rambling way gave vent to a lot of nonsense, which
        Tai-yü did not heed, and Pao-yü went on to inquire: "How old she was when
        she came to the capital? what sights and antiquities she saw on the journey?
        what relics and curiosities there were at Yang Chou? what were the local
        customs and the habits of the people?" Tai-yü made no reply; and Pao-yü
        fearing lest she should go to sleep, and get ill, readily set to work to
        beguile her to keep awake. "Ai yah!" he exclaimed, "at Yang Chou, where your
        official residence is, has occurred a remarkable affair;}
}
{
}

\Verse{95}
{
    \zA{}
}
{
    \eA{have you heard about it?" Tai-yü perceiving that he spoke in earnest, that
        his words were correct and his face serious, imagined that what he referred
        to was a true story, and she therefore inquired what it was? Pao-yü upon
        hearing her ask this question, forthwith suppressed a laugh, and, with a
        glib tongue, he began to spin a yarn. "At Yang Chou," he said, "there's a
        hill called the Tai hill; and on this hill stands a cave called the Lin
        Tzu." "This must all be lies," Tai-yü answered sneeringly, "as I've never
        before heard of such a hill." "Under the heavens many are the hills and
        rivers," Pao-yü rejoined, "and how could you know them all? Wait until I've
        done speaking, when you will be free to express your opinion!" "Go on then,"
        Tai-yü suggested, whereupon Pao-yü prosecuted his raillery.}
}
{
}

\Verse{96}
{
    \zA{}
}
{
    \eA{"In this Lin Tzu cave," he said, "there was once upon a time a whole swarm
        of rat-elves. In some year or other and on the seventh day of the twelfth
        moon, an old rat ascended the throne to discuss matters. 'Tomorrow,' he
        argued, 'is the eighth of the twelfth moon, and men in the world will all be
        cooking the congee of the eighth of the twelfth moon. We have now in our
        cave a short supply of fruits of all kinds, and it would be well that we
        should seize this opportunity to steal a few and bring them over.' Drawing a
        mandatory arrow, he handed it to a small rat, full of aptitude, to go
        forward on a tour of inspection.}
}
{
}

\Verse{97}
{
    \zA{}
}
{
    \eA{The young rat on his return reported that he had already concluded his
        search and inquiries in every place and corner, and that in the temple at
        the bottom of the hill alone was the largest stock of fruits and rice. 'How
        many kinds of rice are there?' the old rat ascertained, 'and how many
        species of fruits?' 'Rice and beans,' the young rat rejoined, 'how many
        barns-full there are, I can't remember; but in the way of fruits there are
        five kinds: 1st, red dates; 2nd, chestnuts; 3rd, ground nuts; 4th, water
        caltrops, and 5th, scented taros.' At this report the old rat was so much
        elated that he promptly detailed rats to go forth; and as he drew the
        mandatory arrow, and inquired who would go and steal the rice, a rat readily
        received the order and went off to rob the rice.}
}
{
}

\Verse{98}
{
    \zA{}
}
{
    \eA{Drawing another mandatory arrow, he asked who would go and abstract the
        beans, when once more a rat took over the arrow and started to steal the
        beans; and one by one subsequently received each an arrow and started on his
        errand. There only remained the scented taros, so that picking again a
        mandatory arrow, he ascertained who would go and carry away the taros:
        whereupon a very puny and very delicate rat was heard to assent. 'I would
        like,' he said, 'to go and steal the scented taros.'}
}
{
}

\Verse{99}
{
    \zA{}
}
{
    \eA{The old rat and all the swarm of rats, upon noticing his state, feared that
        he would not be sufficiently expert, and apprehending at the same time that
        he was too weakly and too devoid of energy, they one and all would not allow
        him to proceed. 'Though I be young in years and though my frame be
        delicate,' the wee rat expostulated, 'my devices are unlimited, my talk is
        glib and my designs deep and farseeing; and I feel convinced that, on this
        errand, I shall be more ingenious in pilfering than any of them.' 'How could
        you be more ingenious than they?' the whole company of rats asked.}
}
{
}

\Verse{100}
{
    \zA{}
}
{
    \eA{'I won't,' explained the young rat, 'follow their example, and go straight
        to work and steal, but by simply shaking my body, and transforming myself, I
        shall metamorphose myself into a taro, and roll myself among the heap of
        taros, so that people will not be able to detect me, and to hear me;
        whereupon I shall stealthily, by means of the magic art of dividing my body
        into many, begin the removal, and little by little transfer the whole lot
        away, and will not this be far more ingenious than any direct pilfering or
        forcible abstraction?' After the whole swarm of rats had listened to what he
        had to say, they, with one voice, exclaimed: 'Excellent it is indeed, but
        what is this art of metamorphosis we wonder?}
}
{
}

\Verse{101}
{
    \zA{}
}
{
    \eA{Go forth you may, but first transform yourself and let us see you.' At these
        words the young rat laughed. 'This isn't a hard task!' he observed, 'wait
        till I transform myself.' "Having done speaking, he shook his body and
        shouted out 'transform,' when he was converted into a young girl, most
        beauteous and with a most lovely face. "'You've transformed yourself into
        the wrong thing,' all the rats promptly added deridingly; 'you said that you
        were to become a fruit, and how is it that you've turned into a young lady?'
        "The young rat in its original form rejoined with a sneering smile: 'You all
        lack, I maintain, experience of the world;}
}
{
}

\Verse{102}
{
    \zA{}
}
{
    \eA{what you simply are aware of is that this fruit is the scented taro, but
        have no idea that the young daughter of Mr. Lin, of the salt tax, is, in
        real truth, a genuine scented taro.'" Tai-yü having listened to this story,
        turned herself round and raising herself, she observed laughing, while she
        pushed Pao-yü: "I'll take that mouth of yours and pull it to pieces! Now I
        see that you've been imposing upon me." With these words on her lips, she
        readily gave him a pinch, and Pao-yü hastened to plead for mercy. "My dear
        cousin," he said, "spare me; I won't presume to do it again; and it's when I
        came to perceive this perfume of yours, that I suddenly bethought myself of
        this old story." "You freely indulge in abusing people," Tai-yü added with a
        smile, "and then go on to say that it's an old story."}
}
{
}

\Verse{103}
{
    \zA{}
}
{
    \eA{But hardly had she concluded this remark before they caught sight of
        Pao-ch'ai walk in. "Who has been telling old stories?" she asked with a
        beaming face; "do let me also hear them." Tai-yü pressed her at once into a
        seat. "Just see for yourself who else besides is here!" she smiled; "he goes
        in for profuse abuses and then maintains that it's an old story!" "Is it
        indeed cousin Pao-yü?" Pao-ch'ai remarked. "Well, one can't feel surprised
        at his doing it; for many have ever been the stories stored up in his brain.
        The only pity is that when he should make use of old stories, he invariably
        forgets them! To-day, he can easily enough recall them to mind, but in the
        stanza of the other night on the banana leaves, when he should have
        remembered them, he couldn't after all recollect what really stared him in
        the face!}
}
{
}

\Verse{104}
{
    \zA{}
}
{
    \eA{and while every one else seemed so cool, he was in such a flurry that he
        actually perspired! And yet, at this moment, he happens once again to have a
        memory!" At these words, Tai-yü laughed. "O-mi-to-fu!" she exclaimed. "You
        are indeed my very good cousin! But you've also (to Pao-yü) come across your
        match. And this makes it clear that requital and retribution never fail or
        err." She had just reached this part of her sentence, when in Pao-yü's rooms
        was heard a continuous sound of wrangling; but as what transpired is not yet
        known, the ensuing chapter will explain.}
}
{
}
